% Pursuing Stacks (cotes 134-2 to 134-8): archive cote 134-3, site page 36 (PDF page 37), typescript page 36.
% typescript. Draft compared: ps2.tex:1730-1799.
% AI-transcribed by Claude Opus 5.5 (claude-opus-5-5, adaptive thinking, effort high), Message Batches API, batch ps1 (msgbatch_01Rv5hCaNTtUt4JHx1yqoHPL), 2026-11-12; not yet checked by eye.
\scanpage{134-3/36}
\mspage{36}
having to make a détour by the"basic"modelizer (Cat)\inserted{we started with}. But here it becomes
about time to try and leave the thin air conjecturing, and find some simple
\margin{ch} and concrete \over{\unsure{c}}{ch}aracterization of test categories, \xout{which} or possibly some
reinforcement still of that notion, which will imply stability under the
product operation.

Here it is tempting to use semi-simplicial techniques though, by lack
of independant foundations of homotopy theory in terms of the modelizer
(Cat). Thus we may want to take the map of semi-simplicial sets correspon-%
ding to (T) when passing to nerves, and express that a) this is a fibration
and b) the fibers of this fibration are "contractible"\inserted[typed]{(or "aspheric")}, which together
will imply that we have a weak equivalence in (Ss sets). Or we may follow
the suggestion of Quillen's set-up, working heuristically in (Cat) as though
we actually knew it is a model category, and expressing that the adjunction
morphism in (T), which is a functor between categories, is actually a
"fibration", and that it's fibers are "contractible" namely weakly equivalent
to a one-point category. I\over{\unsure{n}}{n} any case, a minimum amount of technique seems
needed here, to give the necessary clues for pursuing.

14.3.

30. Since Last week when I stopped with my notes, I got involved a bit
with recalling to mind the "Lego-Teichmüller construction game" for descri-%
bing in a conrete, "visual" way the Teichmüller groups of all possible
types and the main relationships between them, which I had first met with
last year. This and other non mathematical occupations left little time
only for my reflections on homotopy theory, \inserted{which I took up} mainly last night and today.
T\unsure{\over{h}{h}}e focus of attention was the "technical point" of getting a handy carac-%
terization of test categories. The situation I feel is beginning to clari-%
fy somewhat. Last thing I did\hand{,} before reading last weeks' notes and getting
back to the typewriter\hand{,} was to get rid of a delusion which I was dragging
along more or less from the beginning of these notes, namely that our
basic modelizer (Cat), which we were using \xout{\unsure{for us}} as \over{\unsure{p}}{g}iving the most
natural definition of (Hot) in our setting, was a "model category" in the
sense of Quillen, more specifically a "closed model category", where
the "weak equivalences" are the homotopy equivalences of course, and where
cofibrations are just monomorphisms (namely functors injective on objects
and on arrows) - fibrations being defined in terms of these by the
\margin{$|$} Serre-Quillen lifting property. Without\inserted[typed]{even} being so demanding, it turns out
still that there is no reasonable structure of a model catégory on (Cat),
having the correct weak equivalences, and such that the standart \inserted{"Kan} inclusi-%
ons\hand{"} of the following two ordered sets
